% vim: set filetype=tex

\documentclass[12pt,a4paper]{article}

\usepackage[utf8]{inputenc}
\usepackage[english]{babel}

\usepackage{libertine}
\renewcommand*\familydefault{\sfdefault}
\usepackage[T1]{fontenc}

\usepackage{color}
\definecolor{fancy}{rgb}{1,0.5,0}
\definecolor{light}{rgb}{0.5,0.5,0.5}
\color{light}

\usepackage{framed}

\usepackage{setspace}

\usepackage[left=2.5cm,right=2.5cm,top=2cm,bottom=2cm]{geometry}

\usepackage{graphicx}

\usepackage{keystroke}

\usepackage{setspace}
\usepackage[colorlinks=true,linkcolor=light]{hyperref}

\usepackage{titlesec}
\titleformat{\section}{\color{fancy}\normalfont\huge\bfseries}{\thesection}{25pt}{}
\titleformat{\subsection}{\color{fancy}\normalfont\large\bfseries}{\thesubsection}{25pt}{}
\titlespacing{\section}{0cm}{1cm}{1cm}
\titlespacing{\subsection}{0cm}{0.5cm}{0.5cm}

\def\numberline#1{}

\begin{document}

% vim: set filetype=tex

\thispagestyle{empty}

\begin{minipage}{0.25\linewidth}
  \begin{center}
    \includegraphics[width=4cm]{images/misc/bsvlogo_clean.pdf}
  \end{center}
\end{minipage}%
\begin{minipage}{0.15\linewidth}
  \hspace{2.5cm}
\end{minipage}%
\begin{minipage}{0.6\linewidth}
    \begin{flushleft}
      {\color{fancy}\normalfont\huge\bfseries Black Sheep Vision} \\
      \singlespacing
      {\color{light}\normalfont\large\bfseries User Guide\\ \normalfont\large For Version 1.0}
  \end{flushleft}
\end{minipage}%


% vim: set filetype=tex

\thispagestyle{empty}
\tableofcontents

% vim: set filetype=tex

\section{\color{fancy}Chapter 1: Overview}
This chapter will give you a quick overview of common components as well as a short description of the software. Each component is explained in detail in their specific section, respectively.

\subsection{At A Glance}
\begin{figure}[h]
  \centering
  \includegraphics[width=12cm]{images/bsv/Overview.png}
  \caption{Overview}
  \label{fig:overview}
\end{figure}

\subsection{Description}
\begin{tabular}{p{0.2\linewidth}p{0.7\linewidth}}
  \color{fancy}Component & \color{fancy}Purpose \\ \hline
  Menubar & You can use the Menubar to create a new workspace, import your raw files, export data you worked on, adjust the settings or view a quick info dialog. \\ \hline
  View Chooser & By using the View Chooser, you are able to efficiently switch between different views. Due to the fact that all those views are representing your data in a unique way, you are able to improve the explanatory power of your data.  \\ \hline
  Groups Panel & The Groups Panel lets you structure your data by classifying and filtering it in the way you specify it, thus improving the vividness of your information. \\ \hline
  Detail View &  As soon as you have arranged your data, the Detail View shows you specific attributes concerning your selection. \\ \hline
  Subspace Panel & By using the Subspace Panel you are not only in control of the current subspace, but also are able to select the method, the effective outlierness is calculated by. \\ \hline
\end{tabular}

% vim: set filetype=tex

\section{\color{fancy}Chapter 2: Getting Started}

\subsection{What You Need}
In order to visualize your data, you have to import it. This can be done by choosing either a valid file in the \textbf{Arff} specification format, or in the \textbf{CSV} format, during the import process. Furthermore you have to provide a valid \textbf{SSD} file, containing information about all available subspaces.

\subsection{System Requirements}
Depending on the size of your data, the system requirements are heavily fluctuating. Therefore you should take the following specifications with a grain of salt. They should at least give you a rough idea of the hardware you will need.

\begin{tabular}{p{0.3\linewidth}p{0.6\linewidth}}
  \color{fancy}Minimum & \color{fancy}Recommended \\ \hline
  1 GB RAM & 3 GB RAM \\ \hline
  Dual Core, 1.5GHz & Dual Core, 2GHz \\ \hline
  Onboard graphics card & Dedicated graphics card \\ \hline
  Keyboard, Touchpad & Keyboard, Mouse\\ \hline
\end{tabular}

\subsection{Starting the Program}
In order to be able to start the program, you have to have at least the Java Runtime Environment 6 installed. This should be the only step you have to do manually. Starting the program is now as easy as double clicking the supplied JAR file.

\subsection{Improving The Performance}
If you want to squeeze the last bit out of your machine, we discovered that by passing some aggressive parameters to the Java Virtual Machine, a performance boost is noticeable. However, we recommend you try the normal procedure first and come back if you really need to. No matter what, you should adjust the following parameters to your specific circumstances:

\begin{framed}
\begin{verbatim}
java -Xmx1024M -Xms1024M -XX:+UseFastAccessorMethods
     -XX:+AggressiveOpts -XX:+UseAdaptiveGCBoundary
     -XX:SurvivorRatio=16 -XX:+UseParallelGC
     -XX:ParallelGCThreads=4 -jar bsv-1.0.2.jar
\end{verbatim}
\end{framed}

% vim: set filetype=tex

\section{\color{fancy}Chapter 3: Basics}

\subsection{Using The Program}
The program's functions can be accessed by mouse and furthermore, to improve the interactivity, there are selected shortcuts, as listed below. In the top left corner you can find buttons to open the dialogs to import and export new data or to show the settings and info dialog. In the middle of the screen you are able to switch to another view by selecting the desired tab. By clicking on the edge of the left side on your screen you toggle the fading effect of the Groups Panel. This behavior is, similar to the Detail View at the right side, and is possible to use in every view. For detailed information to a specific view, see the subsequent chapters.

\begin{center}
  \begin{spacing}{2.1}
  \begin{tabular}{c|l}
    \color{fancy}Shortcut & \color{fancy}Function\\
    \hline\hline
    \keystroke{F1} & Shows the info dialog.\\
    \Ctrl + \keystroke{N} & Opens the dialog to create a new workspace.\\
    \Ctrl + \keystroke{I} & Opens the dialog to import new data.\\
    \Ctrl + \keystroke{E} & Opens the dialog to export the data to \textbf{.csv} or \textbf{.arff}.\\
    \Ctrl + \Shift + \keystroke{S} & Opens the settings dialog.\\
    \Ctrl + \keystroke{1} & Shows the Table View.\\
    \Ctrl + \keystroke{2} & Shows the Indicator Plot.\\
    \Ctrl + \keystroke{3} & Shows the Scatterplot.\\
    \Ctrl + \keystroke{4} & Shows the Histogram.\\
    \Ctrl + \Tab & Switches to the next view.\\
    \Ctrl + \Shift + \Tab & Switches to the previous view.\\
    \Ctrl + \Spacebar & Shows the dialog to switch the active subspace.\\
    \Ctrl + \Shift + \keystroke{F} & Shows the dialog to rename features.\\
    \Ctrl + \Shift + \keystroke{G} & Creates a new group.\\
    \Ctrl + \keystroke{,} & Shows and hides the groups panel.\\
    \Ctrl + \keystroke{.} & Shows and hides the detail view.\\
  \end{tabular}
  \end{spacing}
\end{center}

\subsection{Import}
The workspace dialog is used to load an existing database or to create a new empty one. This dialog is prompted automatically when the program is started. It accepts only \textbf{.bsv} files. If the dialog is given a file without the \textbf{.bsv} extension or a file that does not exist, it will add the proper extension and create a new empty database. In order to switch the current workspace the user can either use the corresponding button in the top left corner or the appropriate shortcut.
As soon as the user opens an empty database or creates a new one, an import dialog will show up. It needs an input file (\textbf{.arff} or \textbf{.csv}) and an algorithm output file (\textbf{.ssd}) in order to initialize the database with values. As soon as the import successfully finishes, the other components of the program will be initialized. The Table View will be shown as default view.

\subsection{Export}
Exporting a range of elements, is as easy as pressing the export button. You have to select a file (\textbf{.arff} or \textbf{.csv}) and are able to optionally generate an \textbf{.ssd} file.

\subsection{Settings}
By pressing the settings button a dialog will appear, in which you are able to change the language, used to localize the program.

\subsection{Views}
Altogether you can choose from four different views and each has its own advantages to explore the imported data. The Table View is useful to get accurate values from a specific element or to sort the elements by a desired feature. Moreover the Table View is very handy to select single elements by ID or after sorting. The second view is the Feature Indicator, which gives you an overview over all features in one subspace and the trend of one element in comparison to other elements. Furthermore the view provides you tools to sort the order of features and to select elements in several features. For visualizing two specific features you can use the Scatterplot. It displays the features in a two dimensional coordinate system and allows you to zoom and shift the view. In addition you can do a selection with a lasso tool, which emphasizes these elements in other views, too. The fourth and last view is the Histogram. It displays the distribution of elements in a certain feature. This is useful to examine which feature values often occur and which ones are very rare, thus maybe indicating an outlier.

\subsection{Detail View}
On the right side you find the Detail View, which gives you more information about the selected elements. You are able to choose among different statistical measurements, like the median, variance and standard deviation. The calculated details are displayed in combination with the involved groups.

\subsection{Groups Panel}
The Groups Panel (Figure \ref{fig:overview}) is located at the left side of the main window. This panel can be shown or hidden by pressing the arrow button to the left of it. When hidden, the free space is automatically taken by the current view which offers better visualization of the data, especially on small monitors and display resolutions. The number of groups is shown in the top left corner of the panel. This number changes when creating or deleting a group. A group can be created by using the button right next to this number. Below you can find the groups themselves, separated by thin lines. Every group has a name shown in bold letters at the top left side. In order to change this name, you must simply double click on it and type in a new one. There are two buttons on the same line, right to the group name. One for activating and deactivating the group and the other for deleting the group.
After this information, all constraints associated with this group are listed. A constraint can be either static or dynamic and is used to restrict elements in a group.
\textbf{A group with no constraints contains all objects.}
Whenever a constraint is added, the objects in the group are intersected with the objects that satisfy the new constraint. The result of this intersection determines the objects that "stay" in the group after taking the new constraint into consideration. All constraints can be activated, deactivated or deleted.

At the bottom part of each group reside the buttons for adding a constraint, changing the color and adding notes. The color of the group can be either static or dynamic, where static means it will be the same for all objects and dynamic that it will be calculated according to the value of a specified feature.

\subsection{Subspaces}
You can find the section for subspaces at the bottom of the main window (Figure \ref{fig:overview}). On the very left side you can open a dialog to edit the name of the available features or outlierness. Next to that you are able to open a dialog to visualize all existing subspaces and to select a new active subspace. In addition, you can filter the subspaces by entering a search query. All words, split by space, have to be within the name of the subspace. In the middle of this section you can see the currently active subspace and on the right side you can switch the effective outlierness calculation.

\subsection{Effective Outlierness}
The effective outlierness is a statistical measure for the outlierness in one subspace. It is calculated from all outlierness values within one subspace. You are currently able to choose among the calculation by average, minimum or maximum. If the subspace changes the value is automatically adjusted. Furthermore it can be visualized in the same way as other features with real values.


% vim: set filetype=tex

\section{\color{fancy}Chapter 4: Table View}

\begin{figure}[h]
  \centering
  \includegraphics[width=16cm]{images/bsv/TableView.png}
  \caption{Table View}
  \label{fig:tableview}
\end{figure}

\subsection{Functionality}
\begin{tabular}{p{0.2\linewidth}p{0.7\linewidth}}
  \color{fancy}Feature & \color{fancy}Usage \\ \hline
  Display & The Table View gives you a simplistic dump of your data, combined with information about specific IDs and their group membership, easily recognizable by their color field. \\ \hline
  Control & On the right side you find various buttons, to control the table's behavior, for example the decimal places shown or the column width. \\ \hline
  Sort & You are able to sort each column, by clicking on its header. \\ \hline
  Rearrange & If you want to rearrange the columns, you are able to manually drag and drop them in the preferred order. \\ \hline
  Select & Selecting specific elements is as simple as marking them in the table and clicking the select button at the right side of the table. Unselecting elements is done analogically. \\ \hline
\end{tabular}

\subsection{Shortcuts}
\begin{center}
  \begin{spacing}{2.5}
  \begin{tabular}{c|l}
    \color{fancy}Shortcut & \color{fancy}Function\\
    \hline\hline
    \Ctrl + \keystroke{S} & Selects the marked elements.\\
    \Ctrl + \keystroke{A} & Selects all elements.\\
    \Ctrl + \keystroke{D} & Deselects the marked elements.\\
    \Ctrl + \keystroke{R} & Deselects all elements.\\
  \end{tabular}
  \end{spacing}
\end{center}

% vim: set filetype=tex

\section{\color{fancy}Chapter 5: Feature Indicator}

\begin{figure}[h]
  \centering
  \includegraphics[width=16cm]{images/bsv/FeatureIndicator.png}
  \caption{Feature Indicator}
  \label{fig:featureindicator}
\end{figure}

\subsection{Functionality}
\begin{tabular}{p{0.2\linewidth}p{0.7\linewidth}}
  \color{fancy}Feature & \color{fancy}Usage \\ \hline
  Display & The Feature Indicator shows all currently available features side by side, with each value of an element connected by a line.  \\ \hline
  Control & On the right side you find various control mechanisms. You are able to adjust the transparency of the view, independent from each group's transparency, export the current view (e.g. as a \textbf{.svg}) or change the visibility of specific features. \\ \hline
  Rearrange &  You are able to rearrange the features by simple drag and drop gestures, within the list of features. Keep in mind to update the view, by clicking on the update button, below the list, after rearranging.\\ \hline
  Autosort &  By using the autosort functionality, the active features are automatically sorted, according to their correlation coefficient. This should give you good arrangement for a first overview, thus it is recommended to use, if your data consists of a large number of features (i.e. ten or more).  \\ \hline
  Select & You are able to select ranges of a specific feature by using your mouse to drag a bar at the feature axis. To discard all selections, you have to click the clear selection button, at the bottom of the control menu. \\ \hline
\end{tabular}

\subsection{Shortcuts}
\begin{center}
  \begin{spacing}{2.5}
  \begin{tabular}{c|l}
    \color{fancy}Shortcut & \color{fancy}Function\\
    \hline\hline
    \Ctrl + \keystroke{R} & Resets the selection.\\
  \end{tabular}
  \end{spacing}
\end{center}

% vim: set filetype=tex

\section{\color{fancy}Chapter 6: Scatterplot}

\begin{figure}[h]
  \centering
  \includegraphics[width=16cm]{images/bsv/Scatterplot.png}
  \caption{Scatterplot}
  \label{fig:scatterplot}
\end{figure}

\subsection{Functionality}
\begin{tabular}{p{0.2\linewidth}p{0.7\linewidth}}
  \color{fancy}Feature & \color{fancy}Usage \\ \hline
  Display & The Scatterplot visualizes values of two features at a time. \\ \hline
  Control & On the right side you find the control menu, which lets you export the current view (e.g. as a \textbf{.svg}), specify the two features to plot or set the mode of the view. \\ \hline
  Modes & There are three different modes: move, zoom, and select. The move mode acts as default and lets you move the plot. The zoom mode lets you zoom in on a specific position. The select mode offers a tool for selecting an area of values. \\ \hline
  Select & You are able to select areas in this view by using the select mode. To discard all selections, you have to click the clear selection button, at the bottom of the control menu. \\ \hline
\end{tabular}

\subsection{Shortcuts}
\begin{center}
  \begin{spacing}{2.5}
  \begin{tabular}{c|l}
    \color{fancy}Shortcut & \color{fancy}Function\\
    \hline\hline
    \Ctrl + \keystroke{J} & Activates the move mode.\\
    \Ctrl + \keystroke{K} & Activates the zoom mode.\\
    \Ctrl + \keystroke{L} & Activates the select mode.\\
    \Ctrl + \keystroke{S} & Activates the select mode.\\
    \Ctrl + \keystroke{R} & Resets the selection.\\
  \end{tabular}
  \end{spacing}
\end{center}

% vim: set filetype=tex

\section{\color{fancy}Chapter 7: Histogram}

\begin{figure}[h]
  \centering
  \includegraphics[width=16cm]{images/bsv/Histogram.png}
  \caption{Histogram}
  \label{fig:histogram}
\end{figure}

\subsection{Functionality}
\begin{tabular}{p{0.2\linewidth}p{0.7\linewidth}}
  \color{fancy}Feature & \color{fancy}Usage \\ \hline
  Display & The Histogram shows the distribution of your data for one feature at a time. \\ \hline
  Control & On the right side you find the control menu, which lets you export the current view (e.g. as a \textbf{.svg}), specify the feature to plot or set the number of bars to plot, in order to better control the visualization of peaks in your data. \\ \hline
  Select & You are able to select the feature, for which the distribution is to be shown, by selecting it in the control menu. \\ \hline
\end{tabular}

\subsection{Shortcuts}
There are no shortcuts defined for this view, in the current version.

% vim: set filetype=tex

\thispagestyle{empty}

\null
\vfill

\centerline{This work by the Black Sheep Vision Developer Team is licensed under a}
\centerline{Creative Commons Attribution-NonCommercial-ShareAlike License.}

\begin{figure}[b]
  \centering
  \includegraphics[width=4cm]{images/misc/by-nc-sa.pdf}
\end{figure}


\end{document}
